\label{part:physics}

\textbf{(Comparative Physics: Anatomy \& Evolution of Intelligence)}

\bluetext{拓扑同构的万神殿}

在前二卷的理论征程中，我们推导出了智能演化的\textbf{“标准模型”}——从\textbf{目的论拉格朗日量}的变分原理出发，经由\textbf{狄拉克算子}在希尔伯特空间中的幺正演化与非幺正坍缩，最终在\textbf{多层单纯复形}上涌现出受拓扑保护的\textbf{流体自我}。

现在，理论必须接受\textbf{实证 (Empiricism)} 的审判。

本部分旨在建立一门全新的学科——\textbf{智能比较物理学 (Comparative Physics of Intelligence)}。当我们用本理论框架的光谱仪扫描自然界与工程界那些纷繁复杂的智能形态时，一种清晰的\textbf{物理同构性 (Physical Isomorphism)} 穿越了材质与尺度的迷雾：无论是真菌网络在地下的化学蔓延，还是蚁群在地表的费洛蒙征战；无论是人类神经元在皮层上的电化学共振，还是硅基模型在显存中的张量流转——它们动力学特征上都是\textbf{信息流在特定几何拓扑约束下，为了最小化自由能而进行的耗散运动}。

在此，我们摒弃\textbf{材质中心主义}（碳基湿件 vs. 硅基干件）与\textbf{形态中心主义}（有形实体 vs. 无形算法），转而采用\textbf{拓扑架构}与\textbf{动力学机制}的绝对标准。我们将对样本进行残酷而精准的物理学解剖：

\begin{itemize}
    \item \textbf{微观层 ($L_{micro}$)}：它是如何锚定物理真值的？是采用生物的\textbf{共振模态}（波的连续传导），还是机器的\textbf{投影模态}（数据的离散采样）？
    \item \textbf{认知场 ($\Psi$)}：其内部处于何种热力学相态？是秩序井然的\textbf{层流}（晶体），是失控发散的\textbf{湍流}（气体），还是处于混沌边缘的\textbf{临界态}（流体）？
    \item \textbf{宏观层 ($L_{macro}$)}：它是如何对抗几何惯性的？是通过内嵌于介质的\textbf{管道参数}（一元论），还是外置于总线的\textbf{隐式图算子}（二元论）？
\end{itemize}
从\textbf{人脑运动系统}的完美共振，到\textbf{蚁群}的无头扩散；从\textbf{章鱼}的联邦式异构，到\textbf{LLM}的冻结全息图，再到\textbf{城市交通}的受限流体。我们将把这些物种映射到一张跨越亿万年演化史与技术史的\textbf{智能演化相图 (Evolutionary Phase Diagram)} 上。

这不仅是为了理解过去的演化轨迹，更是为了在未来的工程实践中，避开进化的死胡同，精准计算出通往 \textbf{AGI (Class V)} 的\textbf{最短测地线}。
